\documentclass[11pt]{article}
\usepackage[margin=1in]{geometry}
\usepackage{amsmath,amssymb,amsthm}
\usepackage[T1]{fontenc}
\usepackage{lmodern}
\usepackage{microtype}
\usepackage{needspace}
\usepackage[hidelinks]{hyperref}
\newtheorem{proposition}{Proposition}
\newtheorem{theorem}[proposition]{Theorem}
\newcommand{\R}{\mathbb R}
\setlength{\parskip}{3pt}
\setlength{\parindent}{1em}
\title{\vspace{-1.5em}Consequences of Minimum Distance\\in Mathematics and Physics}
\author{J. R. Landers}
\date{November 13, 2026}
\begin{document}
\maketitle
\vspace{-1.2em}

A right triangle whose hypotenuse is the shortest possible length is a small object with a large problem: both its legs must be shorter still. The contradiction is elementary, but it asks a serious question. Which familiar constructions survive if geometry has a lower limit? This exposition follows a route from geometry to dynamics to measurement. Elementary geometric restrictions lead to classical lattice waves; a positivity argument bounds their dispersion; and an established astronomical strategy turns that bound into a conditional test. The aim is to connect these ideas while distinguishing what the cutoff forces from what the chosen physics supplies.

\section{What is being cut off?}
The Planck length is the combination
\begin{equation}
\ell_P=\sqrt{\frac{\hbar G}{c^3}}\simeq1.616\times10^{-35}\ \mathrm{m},
\end{equation}
where $\hbar$ is the reduced Planck constant, $G$ is Newton's gravitational constant, and $c$ is the low-energy speed of light. It is a natural scale for quantum-gravity questions, but it is not an experimentally established smallest distance. Moreover, a limit on measurement precision need not be a gap between possible positions; minimal-length proposals implement different mathematical structures \cite{hossenfelder}.

Here we impose the stronger, explicitly conditional hypothesis. Let $(X,d)$ be a metric space of physical spatial positions, with distance function $d$, and suppose there is a constant $k>0$ such that
\begin{equation}\label{gap}
x\ne y\quad\Longrightarrow\quad d(x,y)\geq k.
\end{equation}
One may set $k=\ell_P$, but none of the geometric deductions needs that identification. Condition \eqref{gap} is called uniform discreteness. It specifies neither a lattice nor integer multiples of $k$, and it does not prescribe motion. It also excludes alternatives such as Kempf's sampling-theoretic construction, where continuous fields have a finite density of independent degrees of freedom \cite{kempf}. Throughout, our distances are spatial distances in a chosen description, not Lorentzian spacetime intervals.

\section{The triangles that remain}
To retain a Euclidean notion of shape, consider triples of distinct points whose mutual distances agree with the side lengths of an ordinary Euclidean triangle. Its vertices are physical positions; the continuum of points along its drawn edges is not assumed to belong to $X$. Requiring all those edge points would already contradict \eqref{gap}.

\begin{proposition}\label{tri}
An admissible Euclidean right triangle has hypotenuse $h\geq\sqrt{2}\,k$. If $\theta$ is either acute angle, then
\begin{equation}\label{angles}
\arcsin(k/h)\leq\theta\leq\frac{\pi}{2}-\arcsin(k/h).
\end{equation}
\end{proposition}
\begin{proof}
Write its legs as $a=h\sin\theta$ and $b=h\cos\theta$. The cutoff requires $a,b\geq k$, so
\begin{equation}
h^2=a^2+b^2\geq2k^2.
\end{equation}
The inequalities $\sin\theta\geq k/h$ and $\cos\theta\geq k/h$ give \eqref{angles}.
\end{proof}
In particular, no such right triangle has hypotenuse $k$. At the threshold $h=\sqrt{2}k$, only the isosceles right shape is permitted by the side cutoff. These are necessary conditions: a particular space $X$ may fail to contain even the permitted triples.

There is a precise way to count the shapes lost, provided we specify a probability measure. Fix $h$ and sample $\theta$ uniformly from $(0,\pi/2)$. The allowed angular interval has length $\pi/2-2\arcsin(k/h)$, whenever it exists. Thus the fraction excluded by the side cutoff alone is
\begin{equation}\label{fraction}
P_{\mathrm{exc}}(h,k)=
\begin{cases}
1,&0<h<\sqrt{2}k,\\[2pt]
\displaystyle\frac{4}{\pi}\arcsin(k/h),&h\geq\sqrt{2}k.
\end{cases}
\end{equation}
The single surviving angle at equality has probability zero. For $h=2k$, two thirds of the sampled shapes are excluded; for $h=10k$, the fraction is approximately $0.128$. Expanding at $k/h=0$ gives
\begin{equation}
P_{\mathrm{exc}}(h,k)=\frac{4k}{\pi h}
+\frac{2k^3}{3\pi h^3}+O\!\left((k/h)^5\right).
\end{equation}
Large triangles therefore recover most shapes under this measure, but arbitrarily thin right triangles remain forbidden. There is no measure-independent ``fraction of all triangles'': choosing sizes and angles differently changes the answer.

\section{Midpoints and motion}
The difficulty with a shortest distance reaches beyond triangles. If a segment of length $L\geq k$ admits $n$ successive rounds of exact Euclidean bisection, its smallest resulting segments have length $L/2^n$. Hence
\begin{equation}\label{bisection}
\frac{L}{2^n}\geq k
\quad\Longrightarrow\quad
n\leq\left\lfloor\log_2(L/k)\right\rfloor.
\end{equation}
This is an upper bound on permitted rounds, not a guarantee that the earlier midpoints exist. Exact, indefinitely repeatable midpoint construction is incompatible with a positive distance gap.

The next observation is a standard connectedness argument, stated here to make the cost of the distance gap explicit.
\begin{proposition}\label{motion}
Under \eqref{gap}, every continuous map $\gamma:[0,T]\to X$, where $T>0$, is constant.
\end{proposition}
\begin{proof}
The open metric ball of radius $k/2$ about any $x\in X$ contains only $x$. Thus every singleton is open, and $X$ has the discrete topology. A continuous image of the connected interval $[0,T]$ is connected, whereas a discrete space has no connected subsets with more than one point.
\end{proof}
Consequently, ordinary continuous point trajectories cannot describe nontrivial motion in this geometry. This does not preclude evolving fields or quantum amplitudes on discrete positions: their values may change continuously with time even though an individual position-valued continuous path cannot. That distinction provides a route from a mathematical obstruction to a dynamical model.

\section{From differences to dispersion}
Choose the particular space $X=k\mathbb Z$, so its positions are $x_n=nk$ for integers $n$. Let $u_n(t)$ be a real scalar wave amplitude at site $n$ and continuous time $t$. Suppose neighboring sites interact according to
\begin{equation}\label{wave}
\ddot u_n=c^2\frac{u_{n+1}-2u_n+u_{n-1}}{k^2}.
\end{equation}
Dots denote time derivatives. This is the classical harmonic-chain wave equation; general finite-range chains and their dispersive behavior are studied, for example, by Mielke and Patz \cite{mielke}. The lattice and evolution rule are additional assumptions, not consequences of \eqref{gap}. Their long-wavelength limit is the usual wave equation.

Insert a plane wave $u_n(t)=\operatorname{Re}\{A e^{i(qnk-\omega t)}\}$, where $A$ is a complex amplitude, $q$ the wave number, and $\omega$ the angular frequency. Equation \eqref{wave} yields
\begin{align}
-\omega^2&=\frac{c^2}{k^2}\bigl(e^{iqk}-2+e^{-iqk}\bigr),\\
\omega^2&=\frac{4c^2}{k^2}\sin^2(qk/2)
=c^2q^2-\frac{c^2k^2q^4}{12}+O(c^2k^4q^6).
\label{dispersion}
\end{align}
The independent wave numbers lie in $-\pi/k\leq q\leq\pi/k$: adding $2\pi/k$ gives the same site values. The continuum derivative has become a finite difference, and its Fourier symbol $-q^2$ has become $-4\sin^2(qk/2)/k^2$. This change in elementary mathematics is what produces dispersion.

For the positive-frequency, right-moving branch $0<q<\pi/k$, a narrow wave packet has group velocity
\begin{equation}\label{velocity}
v_g=\frac{d\omega}{dq}=c\cos(qk/2)
=c\left[1-\frac{k^2q^2}{8}+O((qk)^4)\right].
\end{equation}
Shorter wavelengths propagate more slowly. The usual speed $c$ reappears as $qk\to0$. A spatial lattice selects a preferred rest frame; extending it to three dimensions also generally introduces directional dependence. This scalar example is an illustrative propagation model, not a complete electromagnetic or relativistic theory. Before considering a measurement, we ask whether the slowdown is merely an artifact of coupling each site to its immediate neighbors.

\section{A slowdown shared by a class of models}
Keep the positions $x_n=nk$, but allow symmetric interactions out to $R$ lattice steps:
\begin{equation}\label{generalwave}
\ddot u_n=\sum_{r=1}^{R}g_r
\bigl(u_{n+r}-2u_n+u_{n-r}\bigr),\qquad g_r\geq0.
\end{equation}
Here $R$ is a finite positive integer, $r$ counts lattice steps, and $g_r$ is a coupling with units of inverse time squared. Assume at least one $g_r$ is positive. These are linear, translation-invariant interactions, with equal couplings to the left and right. Nonnegative couplings describe ordinary restoring interactions.

The same operator appears in numerical differentiation. There, positivity obstructs cancellation of the fourth-degree term: Davydov and Schaback prove this for much more general point configurations \cite[\S4.3, Proposition 17]{davydov}. Below we express that familiar obstruction as a sharp bound on the leading wave-speed correction, and also bound the group speed at every differentiable wave number.

For \eqref{generalwave}, define the long-wavelength speed $c$ and effective squared length $\Lambda^2$ by
\begin{equation}\label{moments}
c^2=k^2\sum_{r=1}^{R}g_r r^2,
\qquad
\Lambda^2=k^2\frac{\sum_{r=1}^{R}g_r r^4}{\sum_{r=1}^{R}g_r r^2}.
\end{equation}
\Needspace{9\baselineskip}
\begin{theorem}\label{universal}
For the interactions in \eqref{generalwave}, with $c$ and $\Lambda$ defined by \eqref{moments}, the following bounds hold.
The positive-frequency dispersion branch has group speed $|v_g|\leq c$ wherever differentiable. As $q\to0^+$,
\begin{equation}\label{universalvelocity}
v_g(q)=c\left[1-\frac{\Lambda^2q^2}{8}+O(q^4)\right],
\qquad \Lambda^2\geq k^2.
\end{equation}
Equality $\Lambda^2=k^2$ holds exactly when all couplings beyond nearest neighbors vanish.
\end{theorem}
\begin{proof}
The same plane-wave substitution gives
\begin{align}
\omega^2(q)&=4\sum_{r=1}^{R}g_r\sin^2(rqk/2)\label{generaldispersion}\\
&=c^2q^2\left[1-\frac{\Lambda^2q^2}{12}+O(q^4)\right].
\end{align}
All frequencies are real because the summands in \eqref{generaldispersion} are nonnegative. Taking the positive square root for $q>0$ yields $\omega=cq[1-\Lambda^2q^2/24+O(q^4)]$; differentiating gives \eqref{universalvelocity}. The expansion is for each fixed finite-range model.

To bound its coefficient, set $w_r=g_r r^2/\sum_s g_s s^2$. These weights are nonnegative and sum to one, so
\begin{equation}\label{momentbound}
\Lambda^2=k^2\sum_{r=1}^{R}w_r r^2\geq k^2.
\end{equation}
Equality requires every nonzero weight to have $r=1$. Finally, let the vector $\mathbf f(q)$ have components $f_r(q)=2\sqrt{g_r}\sin(rqk/2)$. Since $\omega=\|\mathbf f\|$, Cauchy--Schwarz gives, wherever the norm is differentiable,
\begin{equation}
|v_g|\leq\|\mathbf f'(q)\|
=\sqrt{\sum_{r=1}^{R}g_r r^2k^2\cos^2(rqk/2)}\leq c.\qedhere
\end{equation}
\end{proof}
The coefficient has a simple meaning: $\Lambda^2$ is a weighted average of squared interaction distances. Every such distance is at least $k$, so the leading slowdown cannot be smaller than in the nearest-neighbor model. Longer-range positive interactions cannot cancel it. The speed bound concerns group velocity; it does not assert a strict causal front for the lattice evolution.

The positivity assumption is essential, even if we retain real frequencies. The classical fourth-order five-point second-derivative formula \cite[Eq.~3.4.23]{nist} supplies a counterexample: allow a negative second-neighbor coupling and choose
\begin{equation}
g_1=\frac{4c^2}{3k^2},\qquad
g_2=-\frac{c^2}{12k^2},\qquad g_r=0\ (r>2).
\end{equation}
Then $k^2(g_1+4g_2)=c^2$ but $g_1+16g_2=0$, so the $q^4$ term in $\omega^2$ cancels. Nevertheless, with $s=\sin(qk/2)$,
\begin{equation}
\omega^2(q)=\frac{4c^2}{k^2}s^2\left(1+\frac{s^2}{3}\right)\geq0.
\end{equation}
Thus the nonzero-frequency modes remain oscillatory, yet the quadratic correction to group velocity disappears. A standard device for improving numerical accuracy becomes, when interpreted as physical dynamics, an example of how the predicted signal can change. The cutoff alone does not determine it.

\section{An arrival-time experiment}
Suppose photon wave packets inherit the lattice dispersion along their propagation direction, with energy $E=\hbar\omega$. This is the extra phenomenological hypothesis required to connect the model to light. In the nearest-neighbor case, define the energy scale $E_k=\hbar c/k$; equation \eqref{velocity} becomes exactly
\begin{equation}
v_g(E)=c\sqrt{1-\frac{E^2}{4E_k^2}}
=c\left[1-\frac{E^2}{8E_k^2}+O\!\left((E/E_k)^4\right)\right].
\end{equation}
For two simultaneously emitted packets of energies $E_h>E_l$, traveling a fixed distance $L$ in a static setting, the leading arrival-time difference is
\begin{equation}\label{delay}
\Delta t=\frac{L}{v_g(E_h)}-\frac{L}{v_g(E_l)}
\simeq\frac{L}{8c}\frac{E_h^2-E_l^2}{E_k^2}
=\frac{Lk^2(E_h^2-E_l^2)}{8\hbar^2c^3}.
\end{equation}
For the full positive-coupling class, $q=E/(\hbar c)+O(E^3)$ at low energy, so Theorem~\ref{universal} replaces $k^2$ in the leading delay by $\Lambda^2$:
\begin{equation}\label{classdelay}
\Delta t\simeq\frac{L\Lambda^2(E_h^2-E_l^2)}{8\hbar^2c^3},
\qquad \Lambda^2\geq k^2.
\end{equation}
The inequality bounds the leading coefficient, not every finite-energy delay. It is sharp: nearest-neighbor interactions attain it. The prediction is therefore robust across this class even without knowing the individual couplings.

The observable signature is a delay quadratic in energy and increasing with propagation distance. Amelino-Camelia and collaborators proposed gamma-ray bursts as probes of energy-dependent light propagation in 1998 \cite{amelino}; their short features, energetic photons, and long paths make small speed differences accumulate. Fermi observations have since constrained both linear and quadratic energy dependence \cite{vasileiou}. A practical analysis compares energy-resolved light curves across sources and distances, fitting propagation delays together with intrinsic emission lags and instrumental timing effects. One delayed energetic photon is insufficient to establish the effect.

For cosmological sources, $L/c$ in \eqref{delay} must be replaced by an expansion-weighted integral. If the same local law holds with constant physical $k$ in the cosmological rest frame and photons undergo the usual leading-order redshift, the prediction is
\begin{equation}\label{cosmo}
\Delta t_{\mathrm{prop}}\simeq
\frac{E_{h,0}^2-E_{l,0}^2}{8E_k^2}
\int_0^{z_s}\frac{(1+z)^2}{H(z)}\,dz.
\end{equation}
Here $z_s$ is source redshift, $H(z)$ the cosmic expansion rate, and the subscript $0$ denotes observed energies. The integral accounts for the higher photon energies earlier in the journey. For the full class, replace $E_k$ by $E_\Lambda=\hbar c/\Lambda$, assuming $\Lambda$ is constant in that frame. This cosmological embedding is another assumption, rather than a result about the static lattice itself.

The scale of the prediction matters. For $k=\ell_P$, one has $E_k\simeq1.22\times10^{28}\ \mathrm{eV}$. Taking $E_h=10^{12}\ \mathrm{eV}$, negligible $E_l$, and $L/c=10^9$ years in the static estimate gives
\begin{equation}
\Delta t\simeq\frac{3.16\times10^{16}\ \mathrm{s}}{8}
\left(\frac{10^{12}}{1.22\times10^{28}}\right)^2
\simeq2.7\times10^{-17}\ \mathrm{s}.
\end{equation}
Such a delay is far below the temporal resolving power of an astronomical burst measurement. The nearest-neighbor model therefore illustrates a testable dependence without promising present sensitivity to Planck-length spacing. The theorem gives a minimum leading coefficient, not a universal upper bound: longer interaction distances can make $\Lambda$ larger than $k$. In the low-energy regime, if an analysis limits the propagation delay to $\tau>0$, \eqref{classdelay} gives the class-wide conditional bound
\begin{equation}
k\lesssim\sqrt{\frac{8\hbar^2c^3\tau}{L(E_h^2-E_l^2)}}.
\end{equation}

A minimum distance restricts shapes and subdivision, and rules out ordinary continuous point trajectories. Evolving wave amplitudes remain possible; on a lattice with positive, symmetric, finite-range interactions, their group speed is bounded and their leading low-energy slowdown has a sharp minimum coefficient. Arrival-time measurements can constrain that entire class, subject to the photon and cosmological assumptions. The negative-coupling example marks the boundary: geometry supplies restrictions, while dynamics decides which restrictions become observable.

\begin{thebibliography}{9}
\small
\setlength{\itemsep}{2pt}
\setlength{\parskip}{0pt}
\bibitem{hossenfelder} S. Hossenfelder, Minimal length scale scenarios for quantum gravity, \emph{Living Reviews in Relativity} \textbf{16} (2013), 2. \href{https://doi.org/10.12942/lrr-2013-2}{doi:10.12942/lrr-2013-2}; \href{https://arxiv.org/abs/1203.6191}{arXiv:1203.6191}.
\bibitem{kempf} A. Kempf, A generalized Shannon sampling theorem, fields at the Planck scale as bandlimited signals, \emph{Physical Review Letters} \textbf{85} (2000), 2873--2876. \href{https://doi.org/10.1103/PhysRevLett.85.2873}{doi:10.1103/PhysRevLett.85.2873}.
\bibitem{mielke} A. Mielke and C. Patz, Uniform asymptotic expansions for the infinite harmonic chain, \emph{Zeitschrift f\"ur Analysis und ihre Anwendungen} \textbf{36} (2017), 437--475. \href{https://doi.org/10.4171/ZAA/1596}{doi:10.4171/ZAA/1596}.
\bibitem{davydov} O. Davydov and R. Schaback, Minimal numerical differentiation formulas, \emph{Numerische Mathematik} \textbf{140} (2018), 555--592. \href{https://doi.org/10.1007/s00211-018-0973-3}{doi:10.1007/s00211-018-0973-3}.
\bibitem{nist} NIST Digital Library of Mathematical Functions, \S3.4, Differentiation, Eq.~3.4.23. \url{https://dlmf.nist.gov/3.4.E23}.
\bibitem{amelino} G. Amelino-Camelia, J. Ellis, N. E. Mavromatos, D. V. Nanopoulos, and S. Sarkar, Tests of quantum gravity from observations of gamma-ray bursts, \emph{Nature} \textbf{393} (1998), 763--765. \href{https://doi.org/10.1038/31647}{doi:10.1038/31647}.
\bibitem{vasileiou} V. Vasileiou et al., Constraints on Lorentz invariance violation from Fermi-Large Area Telescope observations of gamma-ray bursts, \emph{Physical Review D} \textbf{87} (2013), 122001. \href{https://doi.org/10.1103/PhysRevD.87.122001}{doi:10.1103/PhysRevD.87.122001}; \href{https://arxiv.org/abs/1305.3463}{arXiv:1305.3463}.
\end{thebibliography}
\end{document}
